\documentclass[12pt]{article}
\usepackage[utf8]{inputenc}
\usepackage[T1]{fontenc}
\usepackage{hyperref}
\title{Tswa Language Resources for K-12 Curriculum Development}
\author{Compiled by Claude Sonnet 5 (claude-sonnet-5) for Andr\'as Kornai}
\date{}
\begin{document}
\maketitle

\section{Introduction}
Tswa (Xitswa; ISO 639-3 tsc; Glottocode tswa1255) is a Tswa--Ronga Bantu language spoken mainly west of Inhambane in southern Mozambique. Some speakers live in southern Zimbabwe (the Hlengwe dialect) and in South Africa. Many Mozambicans treat it as a dialect of Tsonga, which complicates speaker counts. Joshua Project lists about 1.42 million Tswa in Mozambique, 27,000 in South Africa and 4,400 Hlengwe speakers in Zimbabwe. It uses a Latin alphabet adapted from Tsonga orthography.

\section{Primary Educational Resources}
No Tswa-specific textbooks or school programmes were found in the searches.

\section{Linguistic and Reference Materials}
\begin{itemize}
\item Tswa language (Wikipedia): \url{https://en.wikipedia.org/wiki/Tswa_language}
\item Tswa--Ronga languages (Wikipedia): \url{https://en.wikipedia.org/wiki/Tswa%E2%80%93Ronga_languages}
\item Omniglot page on Tswa: \url{https://omniglot.com/writing/tswa.htm}
\end{itemize}

\section{Religious and Community Resources}
\begin{itemize}
\item Joshua Project language page: \url{https://joshuaproject.net/languages/tsc}
\end{itemize}
Any Tswa scripture would have to be confirmed separately. None was identified in these searches.

\section{Digital and Audio Resources}
Global Recordings Network lists recordings for two Tswa varieties:
\begin{itemize}
\item Tswa: Ndxhonge: \url{https://www.globalrecordings.net/en/language/17659}
\item Tswa: Mandla: \url{https://www.globalrecordings.net/en/language/17658}
\end{itemize}

\section{Curriculum Development Context}
Tswa has a large speaker base but little documentation, and it is closely related to Tsonga and Ronga. Tsonga materials may be a partial resource, but the intelligibility is only partial, so Tswa-specific material would still be needed.
\end{document}